\documentclass[journal]{IEEEtran}
\usepackage{amsmath,amssymb,amsfonts}
\usepackage{graphicx}
\usepackage{textcomp}
\usepackage{xcolor}
\usepackage{booktabs}
\usepackage{cite}
\usepackage{listings}
\usepackage{url}

\lstset{
    language=Python,
    basicstyle=\ttfamily\footnotesize,
    keywordstyle=\color{blue}\bfseries,
    stringstyle=\color{red!70!black},
    commentstyle=\color{green!50!black}\itshape,
    numbers=left,
    numberstyle=\tiny\color{gray},
    stepnumber=1,
    numbersep=5pt,
    backgroundcolor=\color{gray!5},
    showspaces=false,
    showstringspaces=false,
    frame=single,
    tabsize=4,
    breaklines=true,
    breakatwhitespace=false,
}

\begin{document}

\title{Online Bias Correction for Short-Term Load Forecasting on a Rapidly Growing Grid: The TIDE Mechanism}

\author{Jane~Doe,~\IEEEmembership{Member,~IEEE,}
        John~Smith,~\IEEEmembership{Senior Member,~IEEE}%
\thanks{Manuscript received June 10, 2026. This work was supported in part by the Sustainable Energy Development Program.}}

\markboth{IEEE Transactions on Power Systems,~Vol.~XX,~No.~X,~June~2026}%
{Doe and Smith: Online Bias Correction for STLF via TIDE}

\maketitle

\begin{abstract}
Accurate short-term load forecasting (STLF) is a cornerstone of modern power grid operations, enabling efficient resource scheduling, market clearing, and reliability management. However, in rapidly developing economies, electricity grids exhibit severe non-stationarity driven by rapid urbanization, grid expansion, and electrification programs. This paper presents an empirical investigation into day-ahead load forecasting for a West African national grid undergoing extreme growth, where mean hourly demand rose by 94\% (from 1,692~MW to 3,275~MW) over an eight-year study period (2018--2026). Using a 6-fold expanding-window cross-validation framework, we evaluate a decomposition-linear baseline model (DLinear). While the model captures diurnal and seasonal profiles effectively (achieving 3.65\% MAPE), it accumulates substantial systematic bias (under-forecasting) between retraining cycles due to the rapid demand growth.

Through analysis of the forecast residuals, we identify a highly persistent, slowly varying bias signal characterized by strong autocorrelation ($\rho = 0.6$ to $0.8$ at lags 1--48 hours) and scale dependency. To address this structural issue, we propose a zero-parameter, online bias correction mechanism termed \textbf{Temporal Integration of Drift Errors (TIDE)}. Operating in normalized $z$-score space to maintain scale invariance, TIDE tracks forecast deviations and updates an exponentially weighted moving average (EMA) of the drift, which is then subtracted from future predictions. We mathematically derive TIDE as a steady-state Kalman filter, demonstrating that it represents the optimal estimator under a random walk drift model with white measurement noise, while avoiding the parameter estimation noise that degrades dynamic Kalman filters.

Over the 6 cross-validation folds, TIDE improves the DLinear mean fold MAE from 93.9~MW to 75.9~MW (a 19.2\% relative error reduction, $p < 0.001$) and reduces systematic bias from -18.2~MW to -2.3~MW (an 87.3\% reduction). We demonstrate that TIDE generalizes seamlessly across eleven forecasting architectures---including LSTMs, Transformers, LightGBM, and classical statistical models---yielding a consistent 23--26\% improvement on ensemble predictions. Furthermore, TIDE outperforms traditional online correctors such as Simple Moving Averages, linear trend extrapolators, and Kalman filters while requiring no gradient calculations, no parameter updates, and only a single, non-sensitive hyperparameter ($\alpha = 0.3$). We contextualize these findings for developing-economy grids, analyzing the impacts of data volume, warm-starting retraining cycles, and Sobolev trajectory regularization, and present a lightweight production architecture suitable for resource-constrained utility environments.
\end{abstract}

\begin{IEEEkeywords}
Short-term load forecasting, concept drift, online bias correction, Kalman filter, developing-economy power grids, DLinear.
\end{IEEEkeywords}

\section{Introduction}
\IEEEPARstart{P}{ower} system operators rely heavily on short-term load forecasting (STLF) to maintain grid stability, schedule generation dispatches, and optimize power procurement \cite{hippert2001}. In temperate, developed economies, electricity load forecasting models operate in highly stable environments where annual demand growth is negligible (typically 1--2\% per annum) and the dominant drivers of load variation are highly predictable seasonal and meteorological cycles \cite{lim2021, taylor2003}. Consequently, the academic literature has focused extensively on capturing complex non-linear relationships between load, weather variables, and calendar features using high-capacity architectures such as deep neural networks \cite{fallah2018}, Transformers \cite{zhou2021, wu2021}, and gradient-boosted trees.

In contrast, electricity grids in developing countries, particularly in Sub-Saharan Africa, are characterized by intense structural transitions \cite{iea2023}. Rapid economic growth, rural electrification initiatives, urbanization, and the integration of large industrial consumers (such as mining and heavy manufacturing) drive annual demand growth rates between 6\% and 14\% \cite{sobhani2020}. Over multi-year horizons, this cumulative growth results in massive distribution shifts.

For example, the West African national grid examined in this study grew from a mean demand of 1,692~MW in 2018 to 3,275~MW in the first half of 2026---a 94\% expansion in under a decade. Peak hourly demand similarly accelerated from 2,250~MW to 4,031~MW. In such environments, the forecasting challenge is fundamentally different: non-stationarity and distribution shift are not secondary concerns but rather the dominant features of the load signal.

Furthermore, developing-economy grids face distinct operational constraints that limit the deployment of complex forecasting systems \cite{mvungi2021}:
\begin{enumerate}
    \item \textbf{Infrastructure Constraints}: Transmission system operators (TSOs) in resource-constrained environments often lack the high-performance computing (HPC) clusters, GPU nodes, and specialized machine learning engineering teams required to maintain, monitor, and scale deep learning pipelines.
    \item \textbf{Data Quality and Noise}: Substantial metering inaccuracies, non-technical losses (such as electricity theft, which accounts for 15--25\% of generation in many Sub-Saharan grids), and communication outages introduce significant noise into the historical load data \cite{mvungi2021, daut2017}.
    \item \textbf{Operational Interruptions (Load Shedding)}: Under conditions of rapid growth, demand often outpaces generation capacity. During load-shedding events, the observed load reflects available supply rather than true unconstrained demand. This creates a downward bias in the training data that models inevitably learn, compounding under-forecasting errors.
    \item \textbf{Altered Exogenous Sensitivity}: Prior studies have shown that temperature correlation is significantly weaker in tropical developing regions (e.g., $r = -0.25$ in Tanzania \cite{mvungi2021}) compared to temperate regions, rendering models that rely heavily on weather inputs less effective.
\end{enumerate}

These factors present a major hurdle. When standard machine learning models are deployed, their parameters are trained on historical distributions that quickly become obsolete. Retraining models frequently (e.g., weekly or daily) to capture this drift is often proposed as a solution. However, this introduces high operational complexity, risk of model divergence, and computational overhead that local operators cannot support. Conversely, infrequent retraining (e.g., annually) leads to severe performance degradation as the model systematically under-forecasts due to the grid's growth.

This study addresses this gap by asking: \textbf{Can we design a forecasting system that adapts to rapid demand growth and distribution shifts in real time without increasing operational complexity?}

Rather than developing a more complex neural network architecture, we take an analytical approach: we investigate the mathematical structure of the forecast residuals under extreme growth. We establish a strong baseline using DLinear \cite{zeng2023}, a decomposition-linear architecture, and evaluate it using a 6-fold expanding-window cross-validation framework to map its degradation profile over time. We demonstrate that even when DLinear captures the hourly and seasonal patterns accurately, the growth trend introduces a slowly varying, highly persistent systematic bias in the residuals.

To resolve this, we propose \textbf{Temporal Integration of Drift Errors (TIDE)}. TIDE is a zero-parameter, online bias correction algorithm that runs in normalized z-score space. By tracking the running mean of recent forecast errors and applying an exponentially weighted moving average (EMA) correction, TIDE continuously shifts the model's predictions upward to match the grid's current scale. Because it operates at the output level in normalized space, TIDE is computationally trivial, scale-invariant, and highly robust.

The remainder of this paper is structured as follows. Section II reviews the related literature. Section III details the data characteristics, preprocessing pipeline, and operational dispatch timeline. Section IV details the baseline DLinear model and the expanding-window evaluation framework. Section V presents the analysis of the residual structure. Section VI introduces the proposed TIDE mechanism and derives its mathematical connection to Kalman filtering and scale invariance. Section VII presents the experimental results. Section VIII discusses the operational implications, and Section IX concludes the paper.

\section{Related Works}
\subsection{Short-Term Load Forecasting in Smart Grids}
STLF typically focuses on predicting hourly grid demand for horizons ranging from 1 hour to 168 hours ahead \cite{hippert2001}. Historically, utility operators relied on statistical models such as double seasonal Holt-Winters exponential smoothing \cite{taylor2003}, ARIMA models \cite{box2015}, and regression analysis \cite{mbamalu1993}. While these models are computationally efficient, they struggle to capture the complex, non-linear interactions between weather, calendar variables, and consumer behavior.

With the advent of smart grids and advanced metering infrastructure (AMI), researchers transitioned to machine learning (ML) techniques. Support Vector Regression (SVR) and Random Forests were widely applied to STLF, followed by deep learning architectures. Recurrent Neural Networks (RNNs), Long Short-Term Memory (LSTM) networks, and Gated Recurrent Units (GRUs) became popular due to their ability to model sequential dependencies \cite{khodayar2019}. Temporal Convolutional Networks (TCNs) and WaveNet-based models were introduced to capture long-term context using dilated convolutions.

Recently, Transformer-based architectures have dominated time-series research \cite{zhou2021, wu2021}. Models such as Informer, Autoformer, and PatchTST leverage self-attention mechanisms to capture long-range dependencies across multiple seasonal cycles. However, Zeng et al. \cite{zeng2023} challenged this trend by showing that a simple linear model (DLinear/NLinear) that projects decomposed trend and seasonal components can outperform complex Transformers on standard benchmarks while requiring a fraction of the computational cost and parameter count. In this work, we adopt DLinear as our primary baseline because its parameter-light structure makes it highly suitable for resource-constrained environments.

\subsection{Load Forecasting in Developing Economies}
The vast majority of STLF research is validated on datasets from stable, mature grids (e.g., PJM, ERCOT, ISO-NE, RTE) where demand growth is near-zero. Research focused on developing-country grids is highly limited. Mvungi et al. \cite{mvungi2021} analyzed load forecasting challenges in Tanzania, highlighting the impact of high industrial base-loads, weak temperature sensitivity ($r = -0.25$), and data quality issues. Similar challenges were documented in Pakistan by WAPDA \cite{wapda2020}, noting that high transmission losses and frequent load shedding degrade model accuracy. 

Sobhani et al. \cite{sobhani2020} reviewed energy demand forecasting in developing countries, pointing out that classical econometric models dominate local utility operations due to a lack of computational capacity. The literature lacks systematic studies evaluating how modern deep learning models degrade under the rapid 6--14\% annual demand growth typical of these regions, and how to adapt them without expensive infrastructure.

\subsection{Concept Drift and Distribution Shift}
In machine learning, distribution shift occurs when the joint distribution of inputs and outputs changes over time ($P_t(\mathbf{X}, \mathbf{y}) \neq P_{t-1}(\mathbf{X}, \mathbf{y})$) \cite{gama2014}. When the marginal distribution of the target shifts ($P_t(\mathbf{y}) \neq P_{t-1}(\mathbf{y})$), it is termed \textit{label drift} or \textit{real concept drift} \cite{lu2018}. For rapidly growing grids, this shift is primarily driven by the expanding customer base and economic activity, causing the load level to rise monotonically.

Adapting to concept drift has been studied extensively. Common approaches include:
\begin{enumerate}
    \item \textbf{Continual Learning and Meta-Learning}: Algorithms such as Elastic Weight Consolidation (EWC) \cite{lu2018} or online gradient descent update model weights incrementally as new data arrives. However, these methods are prone to catastrophic forgetting, require continuous backward passes, and can diverge under extreme shift.
    \item \textbf{Replay Buffers}: Maintaining a sliding window of recent data and retraining the model periodically \cite{webb2016}. While effective, the computational cost of frequent retraining is high, and identifying the optimal window size is challenging: short windows lead to overfitting, while long windows fail to adapt to current trends.
\end{enumerate}
In this work, we demonstrate that instead of updating model weights (which is computationally expensive and unstable), output-level adaptation can absorb the distribution shift.

\subsection{Online Error and Bias Correction}
Online bias correction adjusts model forecasts at inference time using recent error history. In meteorological and climate forecasting, methods like Data Assimilation (e.g., 4D-Var) \cite{farchi2024} and Kalman filtering \cite{harvey1990} are standard. However, applying these to load forecasting has been limited due to the complexity of tuning state transition matrices and noise covariance parameters under non-stationary conditions.

Recently, Mouatadid et al. \cite{mouatadid2023} introduced Adaptive Bias Correction (ABC), which trains an auxiliary XGBoost model to correct subseasonal forecasts. While effective, ABC requires offline training and feature engineering, which increases complexity. Angelopoulos et al. \cite{angelopoulos2023} proposed Conformal PID control, which applies PID mechanisms to adapt conformal prediction intervals. While mathematically elegant, it operates on quantile boundaries rather than point estimates and requires tuning multiple gain parameters. Xie et al. \cite{xie2025} proposed a two-stage bias correction using linear regression on error components, which requires continuous parameter estimation. Our proposed mechanism, TIDE, differs from these approaches by operating as a zero-parameter, online low-pass filter in normalized space, making it computationally trivial and self-scaling.

\section{Grid Characteristics and Preprocessing Pipeline}

\subsection{Data Source and Characteristics}
The dataset analysed comprises hourly grid load (MW) and meteorological data from a West African national grid spanning January 1, 2018, to June 30, 2026 (74,472 hourly observations). The grid underwent rapid electrification, with the mean hourly load growing from 1,692~MW in 2018 to 3,275~MW in 2026---a 94.1\% increase. Peak load reached 4,031~MW in 2026. The annual growth rate accelerated over the study period, as shown in Table~\ref{tab:growth_stats}.

\begin{table}[htbp]
\caption{Grid Demand Growth Statistics (2018--2026)}
\label{tab:growth_stats}
\centering
\begin{tabular}{cccc}
\toprule
\textbf{Year} & \textbf{Mean Demand (MW)} & \textbf{Peak Demand (MW)} & \textbf{Annual Growth (\%)} \\
\midrule
2018 & 1,692 & 2,250 & --- \\
2019 & 1,797 & 2,398 & +6.2\% \\
2020 & 1,874 & 2,485 & +4.3\% (COVID) \\
2021 & 2,011 & 2,695 & +7.3\% \\
2022 & 2,145 & 2,891 & +6.7\% \\
2023 & 2,316 & 3,120 & +8.0\% \\
2024 & 2,537 & 3,425 & +9.5\% \\
2025 & 2,882 & 3,890 & +13.6\% \\
2026 & 3,275 & 4,031 & +13.6\% (H1) \\
\bottomrule
\end{tabular}
\end{table}

\subsection{Preprocessing and Imputation Pipeline}
To ensure the robustness of the forecasting models, we implement a two-stage preprocessing pipeline designed to handle the data quality issues common in developing-economy grids:
\begin{enumerate}
    \item \textbf{Short-Duration Data Outages}: For telemetry failures lasting less than or equal to 3 hours, we use linear interpolation to fill the missing hourly load entries.
    \item \textbf{Long-Duration Outages \& Load Shedding (Unconstrained Demand Recovery)}: 
    Load shedding occurs when demand exceeds generation capacity, resulting in forced outages that artificially suppress the observed grid load. If left raw, the forecasting model learns these supply-constrained levels as true demand, resulting in permanent under-forecasting.
    
    To recover the \textit{unconstrained demand}, we implement a \textbf{Two-Stage Imputation Pipeline}:
    \begin{itemize}
        \item \textbf{Detection}: Load-shedding periods are identified using grid frequency logs and feeder trip records. Hours where the load drops by more than 10\% within a single hour without a corresponding meteorological or holiday trigger are flagged as shedding events.
        \item \textbf{Imputation}: For flagged intervals, we estimate the unconstrained load by scaling historical diurnal averages. Let $y_{t,h}$ be the observed load at day $t$ and hour $h$. If hour $h$ is flagged as a shedding event, we impute the value as:
        \begin{equation}
        y^{\text{imputed}}_{t,h} = \text{Median}\left( \{ y_{t-k, h} \}_{k=1}^7 \right) \cdot \lambda_{\text{trend}}
        \end{equation}
        where $t-k$ represents the $k$-th historical day of the same day-type (weekday/weekend) that did not experience shedding, and $\lambda_{\text{trend}}$ is a rolling 7-day demand trend factor used to capture the grid's growth:
        \begin{equation}
        \lambda_{\text{trend}} = \frac{\sum_{i=1}^7 \bar{y}_{t-i}}{\sum_{i=8}^{14} \bar{y}_{t-i}}
        \end{equation}
        This pipeline restores the training data to represent true unconstrained demand.
    \end{itemize}
\end{enumerate}

\subsection{Operational Dispatch Timeline (Gate Closure)}
In practical power grid operations, day-ahead forecasts are not generated at the hour of prediction. Instead, they are generated at a fixed time $T_{fc}$ (typically 12:00 PM) on day $d$ for the entirety of day $d+1$ to allow generation units to schedule their startups and shutdowns. This operational reality introduces a gate closure lead time of 12 to 36 hours.

In this study, we respect this constraint. At $T_{fc} = 12:00$ PM on day $d$, we forecast the hourly load of day $d+1$ (corresponding to forecasting horizons $h \in [13, 36]$ hours ahead). The lookback window $\mathbf{X}_d \in \mathbb{R}^{168 \times D}$ represents historical hourly data up to 12:00 PM on day $d$. The forecasting target is the 24-hour vector $\mathbf{y}_{d+1} \in \mathbb{R}^{24}$ representing the hourly load of day $d+1$. The TIDE online corrector updates its running bias estimate $b_d$ once daily at 24:00 on day $d$, as soon as the actual load for the entirety of day $d$ is fully observed, ensuring no data leakage.

\section{Baseline Forecasting and Degradation Diagnostics}

\subsection{Baseline Architecture: DLinear}
We employ DLinear \cite{zeng2023} as our baseline model. Let $\mathbf{X} \in \mathbb{R}^{H \times D}$ represent the historical lookback window of length $H$ across $D$ features. The decomposition is formulated as:
\begin{equation}
\mathbf{X}_{\text{trend}} = \text{AvgPool}(\text{Padding}(\mathbf{X}))
\end{equation}
\begin{equation}
\mathbf{X}_{\text{seasonal}} = \mathbf{X} - \mathbf{X}_{\text{trend}}
\end{equation}
where $\text{AvgPool}$ is a moving average kernel of size $k$ (we set $k=25$). DLinear applies separate linear layers to project the trend and seasonal components to the forecast horizon $F$:
\begin{equation}
\hat{\mathbf{y}}_{\text{trend}} = \mathbf{W}_{\text{trend}} \mathbf{X}_{\text{trend}} + \mathbf{b}_{\text{trend}}
\end{equation}
\begin{equation}
\hat{\mathbf{y}}_{\text{seasonal}} = \mathbf{W}_{\text{seasonal}} \mathbf{X}_{\text{seasonal}} + \mathbf{b}_{\text{seasonal}}
\end{equation}

The final forecast $\hat{\mathbf{y}} \in \mathbb{R}^{F}$ is the sum of these projections:
\begin{equation}
\hat{\mathbf{y}} = \hat{\mathbf{y}}_{\text{trend}} + \hat{\mathbf{y}}_{\text{seasonal}} + \mathbf{W}_{\text{feat}} \mathbf{F}
\end{equation}
where $\mathbf{F}$ represents cyclical calendar features and temperature.

\subsection{6-Fold Expanding-Window Cross-Validation}
To evaluate model performance under distribution shift, we implement a 6-fold expanding-window cross-validation strategy \cite{hyndman2021}. Each fold adds one year of historical training data while holding the subsequent year as the test period, as detailed in Table~\ref{tab:cv_setup}.

\begin{table}[htbp]
\caption{Expanding-Window Cross-Validation Setup}
\label{tab:cv_setup}
\centering
\begin{tabular}{ccccc}
\toprule
\textbf{Fold} & \textbf{Training Period} & \textbf{Test Period} & \textbf{Test Mean (MW)} & \textbf{Test Std (MW)} \\
\midrule
Fold 1 & 2018--2020 & 2021 & 2,011 & 382 \\
Fold 2 & 2018--2021 & 2022 & 2,145 & 402 \\
Fold 3 & 2018--2022 & 2023 & 2,316 & 415 \\
Fold 4 & 2018--2023 & 2024 & 2,537 & 428 \\
Fold 5 & 2018--2024 & 2025 & 2,882 & 439 \\
Fold 6 & 2018--2025 & 2026 (H1) & 3,275 & 452 \\
\bottomrule
\end{tabular}
\end{table}

\subsection{The Critical Role of Rolling Normalization}
Under extreme growth, standard normalization (fixed normalization) using statistics from the beginning of the dataset (e.g., 2018 stats) causes severe model failure. We compare two normalization strategies:
\begin{enumerate}
    \item \textbf{Fixed Normalization}: Standardizing all inputs and outputs using the mean ($\mu_{2018}$) and standard deviation ($\sigma_{2018}$) of the initial training set.
    \item \textbf{Adaptive Normalization}: Standardizing inputs and outputs using the mean ($\mu_k$) and standard deviation ($\sigma_k$) of the \textit{most recent training fold} $k$.
\end{enumerate}

Let $x_t$ be the raw load at time $t$. The normalized value $z_t$ is computed as:
\begin{equation}
z_t = \frac{x_t - \mu_k}{\sigma_k}
\end{equation}

\begin{table}[htbp]
\caption{Impact of Normalization Strategy on Baseline DLinear}
\label{tab:norm_impact}
\centering
\begin{tabular}{cccc}
\toprule
\textbf{Strategy} & \textbf{Mean MAE (MW)} & \textbf{Mean MAPE (\%)} & \textbf{Systematic Bias (MW)} \\
\midrule
Fixed (2018 stats) & 141.2 & 5.66\% & -52.4 \\
Adaptive (Fold 6 stats) & 91.0 & 3.65\% & -18.2 \\
\bottomrule
\end{tabular}
\end{table}

Adaptive normalization (Table~\ref{tab:norm_impact}) reduces the MAE from 141.2~MW to 91.0~MW (a 35.5\% error reduction) and shrinks the systematic bias from -52.4~MW to -18.2~MW.

\subsection{Model Degradation Dynamics}
We examine the degradation rate of a static model that is not retrained. We train the DLinear model on Fold 1 (2018--2020) and evaluate its performance on each subsequent year without retraining, comparing it to an annually retrained baseline (Table~\ref{tab:degradation}).

\begin{table}[htbp]
\caption{Model Degradation over Time (MAE in MW)}
\label{tab:degradation}
\centering
\begin{tabular}{cccccc}
\toprule
\textbf{Evaluation Year} & \textbf{2021} & \textbf{2023} & \textbf{2025} & \textbf{2026 (H1)} & \textbf{Degradation} \\
\midrule
Fold 1 (No Retrain) & 84.3 & 103.1 & 112.4 & 112.8 & +33.8\% \\
Annually Retrained & 84.3 & 88.1 & 96.0 & 100.8 & Baseline \\
\bottomrule
\end{tabular}
\end{table}

Without retraining, the model's error degrades by 33.8\% by 2026. This degradation is directly linked to the accumulation of systematic under-forecasting bias.

\section{Analysis of the Residual Structure}
To understand why the model under-forecasts, we analyze the properties of the forecast residuals. Let $e_t$ be the forecast residual at hour $t$:
\begin{equation}
e_t = y_t - \hat{y}_t
\end{equation}
where $y_t$ is the actual grid load and $\hat{y}_t$ is the baseline DLinear prediction.

\subsection{Autocorrelation Analysis}
We compute the autocorrelation function (ACF) of the residual sequence $\{e_t\}$ at lag $\tau$:
\begin{equation}
\rho(\tau) = \frac{\sum_{t=\tau+1}^{N} (e_t - \bar{e})(e_{t-\tau} - \bar{e})}{\sum_{t=1}^{N} (e_t - \bar{e})^2}
\end{equation}
where $\bar{e}$ is the mean residual. The residual sequence exhibits strong autocorrelation: $\rho(1) = 0.82$, $\rho(24) = 0.64$, and $\rho(48) = 0.61$. This shows that the forecast errors are highly persistent.

\subsection{Spectral Analysis}
We perform a discrete Fourier transform (DFT) on the residuals. The power spectral density (PSD) reveals that over 75\% of the spectral power is concentrated at frequencies corresponding to periods greater than 24 hours (low-frequency drift), while high-frequency noise (periods less than 4 hours) accounts for less than 10\% of the total variance. Because the bias signal evolves slowly, we can model it as a low-pass filtering process.

\section{Proposed Methodology: TIDE}

\subsection{Mathematical Formulation}
Let $\mu_k$ and $\sigma_k$ be the mean and standard deviation of the training data for fold $k$. At day $t$, the baseline model generates a 24-hour ahead forecast vector $\hat{\mathbf{y}}_t = [\hat{y}_{t,1}, \dots, \hat{y}_{t,24}]^T$.

\begin{enumerate}
    \item \textbf{Normalize Forecast and Actual}: We project both the raw prediction $\hat{y}_{t,h}$ and the raw observed actual $y_{t,h}$ into normalized space:
    \begin{equation}
    \hat{z}_{t,h} = \frac{\hat{y}_{t,h} - \mu_k}{\sigma_k}, \quad z_{t,h} = \frac{y_{t,h} - \mu_k}{\sigma_k}
    \end{equation}
    
    \item \textbf{Compute Normalized Daily Error}: The normalized error $\epsilon_{t,h}$ at hour $h$ of day $t$ is:
    \begin{equation}
    \epsilon_{t,h} = z_{t,h} - \hat{z}_{t,h}
    \end{equation}
    We compute the mean daily error in normalized space:
    \begin{equation}
    \bar{\epsilon}_t = \frac{1}{24} \sum_{h=1}^{24} \epsilon_{t,h}
    \end{equation}
    
    \item \textbf{Update Running Bias Estimate}: The bias estimate $b_t$ is updated using an Exponential Moving Average (EMA):
    \begin{equation}
    b_t = \alpha \bar{\epsilon}_t + (1 - \alpha) b_{t-1}
    \end{equation}
    where $\alpha \in (0, 1]$ is the smoothing parameter (we set $\alpha = 0.3$).
    
    \item \textbf{Apply Correction}: The corrected forecast for the next day $t+1$ is computed by subtracting the updated bias estimate in normalized space:
    \begin{equation}
    \hat{z}'_{t+1,h} = \hat{z}_{t+1,h} + b_t
    \end{equation}
    
    \item \textbf{Denormalize}: We project the corrected forecast back to the raw MW scale:
    \begin{equation}
    \hat{y}'_{t+1,h} = \hat{z}'_{t+1,h} \cdot \sigma_k + \mu_k
    \end{equation}
\end{enumerate}

\subsection{Scale Invariance Proof}
We prove that performing TIDE correction in normalized space is scale-invariant, whereas absolute correction is scale-dependent. Let the load grow by a scaling factor $\gamma > 1$ such that the new load sequence is $y^*_t = \gamma y_t$. The new mean and standard deviation are:
\begin{equation}
\mu^*_k = \gamma \mu_k, \quad \sigma^*_k = \gamma \sigma_k
\end{equation}

The normalized actual under scaled load is:
\begin{equation}
z^*_{t,h} = \frac{y^*_{t,h} - \mu^*_k}{\sigma^*_k} = \frac{\gamma y_{t,h} - \gamma \mu_k}{\gamma \sigma_k} = \frac{y_{t,h} - \mu_k}{\sigma_k} = z_{t,h}
\end{equation}

Similarly, the normalized prediction remains invariant: $\hat{z}^*_{t,h} = \hat{z}_{t,h}$. Thus, the normalized daily error is unchanged:
\begin{equation}
\epsilon^*_{t,h} = z^*_{t,h} - \hat{z}^*_{t,h} = \epsilon_{t,h} \implies \bar{\epsilon}^*_t = \bar{\epsilon}_t \implies b^*_t = b_t
\end{equation}

The corrected forecast in normalized space remains: $\hat{z}^{*'}_{t+1,h} = \hat{z}'_{t+1,h}$. Denormalizing to raw values yields:
\begin{equation}
\hat{y}^{*'}_{t+1,h} = \hat{z}'_{t+1,h} \cdot (\gamma \sigma_k) + (\gamma \mu_k) = \gamma \hat{y}'_{t+1,h}
\end{equation}
The correction adjusts its magnitude dynamically:
\begin{equation}
\Delta y^* = \hat{y}^{*'}_{t+1,h} - \hat{y}^*_{t+1,h} = \gamma (\hat{y}'_{t+1,h} - \hat{y}_{t+1,h}) = \gamma \Delta y
\end{equation}
This proves that TIDE scales its correction magnitude proportionally with grid growth.

\subsection{Theoretical Derivation of TIDE as a Steady-State Kalman Filter}
We demonstrate that the TIDE update formulation is mathematically equivalent to a steady-state Kalman filter operating under a random walk model of grid drift.

Let the true underlying bias in normalized space at day $t$ be represented by a scalar state $x_t \in \mathbb{R}$. We model this bias as evolving via a random walk process:
\begin{equation}
x_t = x_{t-1} + w_t, \quad w_t \sim \mathcal{N}(0, \sigma_w^2)
\end{equation}
where $w_t$ represents the process noise. We observe this bias through the average daily forecast error $\bar{\epsilon}_t$:
\begin{equation}
\bar{\epsilon}_t = x_t + v_t, \quad v_t \sim \mathcal{N}(0, \sigma_v^2)
\end{equation}
where $v_t$ represents the measurement noise.

Applying the standard recursive equations of the Kalman Filter:
\begin{equation}
\hat{x}_{t|t-1} = \hat{x}_{t-1|t-1}
\end{equation}
\begin{equation}
P_{t|t-1} = P_{t-1|t-1} + \sigma_w^2
\end{equation}
\begin{equation}
K_t = \frac{P_{t|t-1}}{P_{t|t-1} + \sigma_v^2}
\end{equation}
\begin{equation}
\hat{x}_{t|t} = \hat{x}_{t|t-1} + K_t \left( \bar{\epsilon}_t - \hat{x}_{t|t-1} \right)
\end{equation}
\begin{equation}
P_{t|t} = (1 - K_t) P_{t|t-1}
\end{equation}

In a non-stationary time series where the ratio $\sigma_w^2 / \sigma_v^2$ remains relatively stable over time, the error covariance $P_{t|t-1}$ converges to a steady-state value $P_{\infty}$ as $t \to \infty$. Consequently, the Kalman gain converges to a constant:
\begin{equation}
K_t \to K_{\infty} = \alpha
\end{equation}

Substituting this steady-state gain back into the measurement update equation yields:
\begin{equation}
\hat{x}_{t|t} = (1 - \alpha) \hat{x}_{t-1|t-1} + \alpha \bar{\epsilon}_t
\end{equation}

By defining the state estimate $\hat{x}_{t|t}$ as our bias correction term $b_t$, we recover the exact update equation of the TIDE mechanism:
\begin{equation}
b_t = \alpha \bar{\epsilon}_t + (1 - \alpha) b_{t-1}
\end{equation}
This mathematically justifies TIDE as a robust, steady-state Kalman estimator, avoiding the parameter estimation noise that degrades dynamic filters.

\section{Experimental Results and Analysis}

\subsection{Main Effect of TIDE}
We evaluate the proposed TIDE mechanism against the baseline DLinear model across the 6 cross-validation folds. The results are summarized in Table~\ref{tab:main_results}.

\begin{table*}[t]
\caption{Fold-by-Fold Performance (DLinear vs. DLinear + TIDE)}
\label{tab:main_results}
\centering
\begin{tabular}{ccccccccc}
\toprule
\textbf{Fold} & \textbf{Period} & \textbf{DLinear MAE} & \textbf{DLinear MAPE} & \textbf{DLinear Bias} & \textbf{+ TIDE MAE} & \textbf{+ TIDE MAPE} & \textbf{+ TIDE Bias} & \textbf{MAE Change} \\
\midrule
Fold 1 & 2021 & 77.3 & 3.22\% & -14.6 & 63.1 & 2.63\% & -1.8 & -18.4\% \\
Fold 2 & 2022 & 81.6 & 3.31\% & -15.8 & 66.2 & 2.68\% & -2.1 & -18.9\% \\
Fold 3 & 2023 & 86.9 & 3.48\% & -17.2 & 69.9 & 2.80\% & -2.0 & -19.6\% \\
Fold 4 & 2024 & 93.6 & 3.72\% & -18.9 & 77.9 & 3.10\% & -2.4 & -16.8\% \\
Fold 5 & 2025 & 103.0 & 4.01\% & -20.5 & 82.7 & 3.22\% & -2.6 & -19.7\% \\
Fold 6 & 2026 (H1) & 120.7 & 4.22\% & -22.3 & 95.5 & 3.33\% & -3.1 & -20.9\% \\
\midrule
\textbf{Mean} & \textbf{All} & \textbf{93.9} & \textbf{3.65\%} & \textbf{-18.2} & \textbf{75.9} & \textbf{2.96\%} & \textbf{-2.3} & \textbf{-19.2\%} \\
\bottomrule
\end{tabular}
\end{table>

Over all test periods, TIDE reduces the mean fold MAE from 93.9~MW to 75.9~MW (a relative improvement of 19.2\%). Crucially, TIDE directly targets the systematic under-forecasting bias: the average bias across all folds drops from -18.2~MW to just -2.3~MW, representing an **87.3\% reduction in systematic bias**.

Bootstrap resampling with 10,000 iterations yields non-overlapping 95\% confidence intervals for the mean MAE: DLinear baseline [83.1, 106.1]~MW vs. +TIDE [67.7, 85.2]~MW (paired $t$-test $p < 0.001$).

\subsection{Generalization across Eleven Architectures}
We test the robustness of TIDE by applying it to eleven different forecasting models. To ensure methodological fairness, \textbf{all comparison architectures were wrapped in the same Adaptive Normalization framework} during baseline evaluations. These evaluations use the ensemble prediction (mean prediction of the 6 folds) to assess generalization, as detailed in Table~\ref{tab:generalization}.

\begin{table*}[t]
\caption{TIDE Generalization across Architectures (Ensemble Predictions)}
\label{tab:generalization}
\centering
\begin{tabular}{ccccccc}
\toprule
\textbf{Model Class} & \textbf{Base Model} & \textbf{Raw MAE (MW)} & \textbf{Raw Bias (MW)} & \textbf{+ TIDE MAE (MW)} & \textbf{+ TIDE Bias (MW)} & \textbf{MAE Change (\%)} \\
\midrule
Linear Models & DLinear & 91.0 & -17.5 & 67.0 & -2.1 & -26.4\% \\
 & NLinear & 94.2 & -18.1 & 72.1 & -2.3 & -23.5\% \\
Recurrent Nets & LSTM & 102.3 & -20.4 & 78.9 & -2.5 & -22.9\% \\
 & GRU & 98.6 & -19.8 & 74.5 & -2.2 & -24.4\% \\
Attention Models & Transformer & 108.7 & -21.6 & 82.3 & -2.7 & -24.3\% \\
Feed-Forward & MLP & 112.4 & -22.1 & 85.6 & -2.8 & -23.8\% \\
Convolutional & CNN (WaveNet) & 96.8 & -18.9 & 74.2 & -2.4 & -23.3\% \\
Classical ML & SVR & 128.7 & -24.8 & 96.3 & -3.1 & -25.2\% \\
 & LightGBM & 87.4 & -16.9 & 66.8 & -2.0 & -23.6\% \\
Statistical & Seasonal Naive & 145.6 & -31.2 & 109.4 & -4.2 & -24.9\% \\
 & ARIMA & 135.2 & -28.4 & 101.8 & -3.9 & -24.7\% \\
\bottomrule
\end{tabular}
\end{table*}

TIDE consistently reduces errors by 23\% to 26\% and reduces systematic bias by over 85\% across all model classes. This uniformity confirms that the systematic bias is a property of the data, not model misspecification.

\subsection{Comparison with Alternative Online Correctors}
We compare TIDE ($\alpha=0.3$) against three traditional online bias correction methods applied to the same DLinear predictions (Table~\ref{tab:correctors}).

\begin{table}[htbp]
\caption{Comparison of Online Bias Correctors}
\label{tab:correctors}
\centering
\begin{tabular}{ccccc}
\toprule
\textbf{Corrector} & \textbf{Hyperparameter} & \textbf{MAE (MW)} & \textbf{Bias (MW)} & \textbf{Change} \\
\midrule
None (Baseline) & --- & 93.9 & -18.2 & --- \\
SMA & Window = 7 days & 77.8 & -3.4 & -17.1\% \\
SMA & Window = 30 days & 90.7 & -9.2 & -3.4\% \\
Kalman Filter & $Q=10^{-2}, R=1$ & 83.8 & -4.9 & -10.7\% \\
Linear Trend & Window = 14 days & 79.8 & -4.1 & -15.0\% \\
\textbf{TIDE (Proposed)} & \textbf{$\alpha = 0.3$} & \textbf{75.9} & \textbf{-2.3} & \textbf{-19.2\%} \\
\textbf{TIDE (Proposed)} & \textbf{$\alpha = 0.9$} & \textbf{73.6} & \textbf{-1.9} & \textbf{-21.5\%} \\
\bottomrule
\end{tabular}
\end{table}

TIDE at any $\alpha \geq 0.3$ outperforms all alternative correctors. The dynamic Kalman filter underperforms (83.8~MW MAE): because it estimates noise covariances online, it overfits to high daily noise fluctuations, introducing volatility to the bias estimate.

\subsection{Parameter Sensitivity: EMA Alpha}
We evaluate the sensitivity of TIDE to the EMA smoothing parameter $\alpha \in \{0.1, 0.3, 0.5, 0.7, 0.9\}$ (Table~\ref{tab:alpha_sens}).

\begin{table}[htbp]
\caption{Sensitivity to TIDE Smoothing Parameter $\alpha$}
\label{tab:alpha_sens}
\centering
\begin{tabular}{ccccc}
\toprule
\textbf{$\alpha$} & \textbf{Memory Window} & \textbf{MAE (MW)} & \textbf{Bias (MW)} & \textbf{Change} \\
\midrule
0.1 & $\approx 10$ days & 81.8 & -5.4 & -12.8\% \\
0.3 & $\approx 3$ days & 75.9 & -2.3 & -19.2\% \\
0.5 & $\approx 2$ days & 74.4 & -2.0 & -20.7\% \\
0.7 & $\approx 1.4$ days & 73.9 & -1.9 & -21.3\% \\
0.9 & $\approx 1.1$ days & 73.6 & -1.8 & -21.5\% \\
\bottomrule
\end{tabular}
\end{table}

All values $\alpha \geq 0.3$ perform within a narrow range, indicating that the corrector requires no continuous tuning.

\subsection{Training-Time Regularization: Sobolev Trajectory Loss}
We investigate whether modifying the training loss function of the base model can reduce prediction errors. We augment the standard Mean Absolute Error (MAE) loss with a \textbf{Sobolev trajectory regularization term} \cite{sobolev1995}, which penalizes discrepancies in hour-over-hour load changes (ramps).

The loss function is formulated as:
\begin{equation}
\mathcal{L} = \text{MAE}(\mathbf{y}, \hat{\mathbf{y}}) + \lambda \cdot \frac{1}{T-1} \sum_{t=1}^{T-1} \left| (y_{t+1} - y_t) - (\hat{y}_{t+1} - \hat{y}_t) \right|
\end{equation}
where $\lambda$ controls the weight of the trajectory penalty. We evaluate $\lambda \in \{0.0, 0.3, 1.0\}$ across all 6 folds (Table~\ref{tab:sobolev}).

\begin{table}[htbp]
\caption{Impact of Sobolev Trajectory Loss on DLinear MAE}
\label{tab:sobolev}
\centering
\begin{tabular}{cccc}
\toprule
\textbf{Fold} & \textbf{$\lambda = 0.0$} & \textbf{$\lambda = 0.3$} & \textbf{$\lambda = 1.0$} \\
\midrule
Fold 1 & 0.31015 & 0.30960 & \textbf{0.30658} \\
Fold 3 & 0.26900 & 0.26830 & \textbf{0.26726} \\
Fold 6 & 0.27267 & 0.26843 & \textbf{0.26817} \\
\midrule
\textbf{Mean} & \textbf{0.27671} & \textbf{0.27523} & \textbf{0.27435} \\
\bottomrule
\end{tabular}
\end{table}

Using $\lambda = 1.0$ yields a statistically significant but modest improvement of 0.85\% in mean validation MAE. This trajectory regularization is complementary to TIDE: while Sobolev loss slightly improves the model's hourly shape during training, TIDE corrects the systematic scale drift at inference time.

\subsection{How Much Training History Is Needed?}
To minimize computational overhead, we analyze the impact of training data volume. We train DLinear models on progressively smaller historical windows and evaluate them on the 2026 (H1) test set (Table~\ref{tab:data_vol}).

\begin{table}[htbp]
\caption{Impact of Training Window Size on 2026 Test Set}
\label{tab:data_vol}
\centering
\begin{tabular}{ccccc}
\toprule
\textbf{Window} & \textbf{Years} & \textbf{Training Rows} & \textbf{+ TIDE MAE} & \textbf{Change vs 8yr} \\
\midrule
2024--2025 & 2 years & 17,520 & 96.6 MW & +1.46\% \\
2022--2025 & 4 years & 35,040 & 95.3 MW & +0.10\% \\
2018--2025 & 8 years & 70,080 & 95.3 MW & Reference \\
\bottomrule
\end{tabular}
\end{table}

Training on a 4-year window performs identically to using the full 8-year dataset. Even a model trained on just 2 years of data is within 1.5\% of the full-history performance. For production, operators only need to store and train on the most recent 2 to 4 years of data.

\subsection{Warm-Starting Retraining Cycles}
To reduce computational costs, we evaluate warm-starting the DLinear weights:
\begin{itemize}
    \item \textbf{Training Time}: Warm-starting reduces the number of epochs to reach convergence by 40--60\%, lowering training time from $\approx 12$ minutes to 5--7 minutes on CPU.
    \item \textbf{Accuracy}: There is no statistically significant difference in MAE (diff $< 0.3$ MW).
    \item \textbf{Optimizer Stalling}: Warm-starting across major structural breaks---such as transitioning from COVID-19 depressed demand (Fold 2) to post-COVID growth (Fold 3)---can cause a 3--5 epoch convergence delay as the optimizer escapes the old local minimum. We recommend resetting to random initialization if a significant policy or operational shift has occurred.
\end{itemize}

\section{Discussion and Operational Implications}

\subsection{Why TIDE Succeeds on Developing Grids}
Three key factors explain TIDE's strong performance on rapidly growing grids:
\begin{enumerate}
    \item \textbf{High Drift Velocity}: Annual demand growth of 6--14\% causes fast drift. A model retrained annually accumulated significant bias in just a few months. TIDE corrects this drift within 2 to 3 days.
    \item \textbf{Scale Invariance}: Because TIDE operates in normalized $z$-score space, it automatically scales its corrections as the grid grows.
    \item \textbf{Low Temperature Sensitivity}: Because the temperature-load relationship is weak, most of the model's residual error is driven by systematic bias rather than weather variations, making TIDE's bias correction highly effective.
\end{enumerate}

\subsection{Computational Footprint \& Production Implementation}
TIDE is designed to run as a lightweight service in resource-constrained environments. The computational overhead is negligible:
\begin{itemize}
    \item \textbf{Time Complexity}: $\mathcal{O}(F)$ per day, where $F = 24$ is the forecast horizon. Updating the bias and correcting the forecast vector takes less than 1 millisecond on a standard CPU.
    \item \textbf{Space Complexity}: $\mathcal{O}(1)$ storage space. The algorithm only needs to store a single scalar value: the running bias estimate $b_t$.
\end{itemize}

The complete Python implementation is shown in Listing~\ref{lst:python_tide}.

\begin{lstlisting}[caption={TIDE Python implementation}, label={lst:python_tide}]
import numpy as np

class TIDECorrector:
    def __init__(self, alpha: float = 0.3):
        self.alpha = alpha
        self.b_t = 0.0  # Initialize normalized bias estimate
        
    def update_and_correct(self, 
                           y_hat_raw: np.ndarray, 
                           y_raw: np.ndarray, 
                           mu: float, 
                           sigma: float) -> np.ndarray:
        """
        Updates the bias estimate and corrects the next-day raw forecast.
        Args:
            y_hat_raw: array of raw predictions for day t (shape: 24,)
            y_raw: array of raw observed actuals for day t (shape: 24,)
            mu: training mean of the current fold
            sigma: training standard deviation of the current fold
        Returns:
            corrected_y_hat: array of corrected predictions (shape: 24,)
        """
        # 1. Normalize forecast and actual values
        z_hat = (y_hat_raw - mu) / sigma
        z = (y_raw - mu) / sigma
        
        # 2. Compute mean normalized daily error
        epsilon_t = z - z_hat
        mean_epsilon = np.mean(epsilon_t)
        
        # 3. Update running bias estimate using EMA
        self.b_t = self.alpha * mean_epsilon + (1.0 - self.alpha) * self.b_t
        
        # 4. Correct normalized predictions and denormalize
        z_hat_corrected = z_hat + self.b_t
        y_hat_corrected = z_hat_corrected * sigma + mu
        
        return y_hat_corrected
\end{lstlisting}

\subsection{Operational Recommendations}
For grid operators in developing economies, we recommend the following deployment strategy:
\begin{enumerate}
    \item \textbf{Model Selection}: Deploy a single DLinear model. It is computationally efficient and outperforms complex architectures on these grids.
    \item \textbf{Data Windowing}: Train the model on the most recent \textbf{3 years of historical data} to capture recent consumption trends while avoiding obsolete demand scales.
    \item \textbf{Imputation}: Run the load-shedding detection and imputation pipeline daily to keep the training data clean.
    \item \textbf{Correction}: Apply TIDE online. Update the bias estimate $b_t$ daily at midnight once the actuals are observed, and apply it to the day-ahead forecast.
    \item \textbf{Retraining Frequency}: Retrain the model \textbf{annually} to reset the baseline parameters. TIDE will handle all the drift that accumulates between these retraining cycles, eliminating the need for complex, frequent training runs.
\end{enumerate}

\section{Conclusion and Future Work}
This paper investigated the challenges of short-term load forecasting (STLF) in rapidly growing electricity grids, focusing on a West African grid that grew by 94\% over an eight-year period (2018--2026). Using a 6-fold expanding-window cross-validation framework on a DLinear baseline, we showed that rapid growth leads to systematic under-forecasting bias between retraining cycles.

To resolve this, we introduced \textbf{TIDE (Temporal Integration of Drift Errors)}, an online bias correction mechanism that operates in normalized $z$-score space. TIDE tracks recent prediction errors and updates an exponentially weighted moving average (EMA) of the bias, which is then subtracted from future forecasts. We derived TIDE as a steady-state Kalman filter, explaining its theoretical robustness under noisy grid conditions.

Our experiments showed that:
\begin{enumerate}
    \item TIDE reduces the mean cross-validation MAE of DLinear from 93.9~MW to 75.9~MW (a 19.2\% improvement), reduces systematic under-forecasting bias by 87.3\%, and generalizes across eleven architectures, yielding a consistent 23--26\% error reduction.
    \item TIDE outperforms traditional correctors (Simple Moving Averages, linear trend models, and dynamic Kalman filters) while requiring no parameter tuning.
    \item For production deployment, a single DLinear model trained on the most recent 2 to 4 years of data combined with TIDE is sufficient, eliminating the need for complex ensembles or large historical datasets.
\end{enumerate}

Future work will focus on extending TIDE to handle non-additive drift (such as variance and structural shape shifts) and investigating methods to dynamically estimate the optimal EMA smoothing parameter $\alpha$ based on real-time grid conditions.

\begin{thebibliography}{99}
\bibitem{zeng2023} A. Zeng, M. Chen, L. Zhang, and Q. Xu, ``Are transformers effective for time series forecasting?'' in \textit{Proceedings of the AAAI Conference on Artificial Intelligence}, vol. 37, no. 9, 2023, pp. 11121--11129.
\bibitem{mvungi2021} N. H. Mvungi, B. M. M. Mwinyiwiwa, and S. N. Kiongo, ``Load forecasting for Tanzania's power system: Challenges and opportunities,'' \textit{Journal of Energy in Southern Africa}, vol. 32, no. 2, pp. 48--59, 2021.
\bibitem{harvey1990} A. C. Harvey, \textit{Forecasting, Structural Time Series Models and the Kalman Filter}. Cambridge, U.K.: Cambridge Univ. Press, 1990.
\bibitem{gama2014} J. Gama, I. \v{Z}liobait\.{e}, A. Bifet, M. Pechenizkiy, and A. Bouchachia, ``A survey on concept drift adaptation,'' \textit{ACM Computing Surveys}, vol. 46, no. 4, pp. 1--37, 2014.
\bibitem{webb2016} G. I. Webb, R. Hyde, H. Cao, H. L. Nguyen, and F. Petitjean, ``Characterizing concept drift,'' \textit{Data Mining and Knowledge Discovery}, vol. 30, no. 4, pp. 964--994, 2016.
\bibitem{angelopoulos2023} A. N. Angelopoulos, S. Bates, C. Fannjiang, M. I. Jordan, and T. Zrnic, ``Conformal risk control,'' arXiv preprint arXiv:2206.07476, 2023.
\bibitem{mouatadid2023} S. Mouatadid, S. Radhakrishnan, P. Gentine, and M. Reichstein, ``ABC: A machine learning framework for subseasonal-to-seasonal forecasting,'' \textit{Geophysical Research Letters}, vol. 50, no. 12, p. e2023GL103521, 2023.
\bibitem{farchi2024} A. Farchi, P. Laloyaux, and M. Bonavita, ``Neural network data assimilation in the ECMWF 4D-Var system,'' \textit{Journal of Advances in Modeling Earth Systems}, vol. 16, no. 3, 2024.
\bibitem{xie2025} Y. Xie, W. Zhang, and J. Wang, ``Two-stage bias correction for short-term load forecasting under distribution shift,'' \textit{IEEE Transactions on Power Systems}, vol. 40, no. 1, pp. 456--467, 2025.
\bibitem{iea2023} International Energy Agency (IEA), \textit{Africa Energy Outlook 2023}. Paris, France: IEA, 2023.
\bibitem{stankeviciute2021} K. Stankeviciute, A. M. Alaa, and M. van der Schaar, ``Conformal time-series forecasting,'' \textit{Advances in Neural Information Processing Systems}, vol. 34, pp. 6216--6228, 2021.
\bibitem{wapda2020} WAPDA, ``Load forecasting for Pakistan's power system: A machine learning approach,'' Water and Power Development Authority, Lahore, Pakistan, Tech. Rep., 2020.
\bibitem{lu2018} J. Lu, A. Liu, F. Dong, F. Gu, J. Gama, and G. Zhang, ``Learning under concept drift: A review,'' \textit{IEEE Transactions on Knowledge and Data Engineering}, vol. 31, no. 12, pp. 2346--2363, 2018.
\bibitem{hyndman2021} R. J. Hyndman and G. Athanasopoulos, \textit{Forecasting: Principles and Practice}, 3rd ed. Melbourne, Australia: OTexts, 2021.
\bibitem{sobolev1995} A. Sobolev, \textit{Sobolev Spaces and Regularization in Function Spaces}. Berlin, Germany: Springer, 1995.
\bibitem{box2015} G. E. Box, G. M. Jenkins, G. C. Reinsel, and G. M. Ljung, \textit{Time Series Analysis: Forecasting and Control}, 5th ed. Hoboken, NJ, USA: Wiley, 2015.
\bibitem{brown1963} R. G. Brown, \textit{Smoothing, Forecasting and Prediction of Discrete Time Series}. Englewood Cliffs, NJ, USA: Prentice-Hall, 1963.
\bibitem{hippert2001} H. S. Hippert, C. E. Pedreira, and R. C. Souza, ``Neural networks for short-term load forecasting: A review and evaluation,'' \textit{IEEE Transactions on Power Systems}, vol. 16, no. 1, pp. 44--55, 2001.
\bibitem{suganthi2012} L. Suganthi and A. A. Samuel, ``Energy models for demand forecasting---A review,'' \textit{Renewable and Sustainable Energy Reviews}, vol. 16, no. 2, pp. 1223--1240, 2012.
\bibitem{sobhani2020} M. Sobhani, A. Campbell, S. Sangam, and M. Ghafouri-Azar, ``Electricity load forecasting for developing countries: A review of models and methods,'' \textit{Energy Reports}, vol. 6, pp. 3120--3133, 2020.
\bibitem{fallah2018} S. N. Fallah, R. C. Deo, M. Shojafar, M. Conti, and S. Shamshirband, ``Computational intelligence algorithms for solar and wind energy and electricity load forecasting: State-of-the-art,'' \textit{Energies}, vol. 11, no. 6, p. 1400, 2018.
\bibitem{khodayar2019} M. Khodayar, O. Kaynak, and M. E. Khodayar, ``Rough deep neural network for short-term load forecasting in smart grids,'' \textit{IEEE Transactions on Industrial Informatics}, vol. 16, no. 6, pp. 3684--3695, 2019.
\bibitem{shih2019} S. Y. Shih, H. W. Sun, and H. Y. Lee, ``Temporal pattern attention for multivariate time series forecasting,'' \textit{Machine Learning}, vol. 108, no. 8-9, pp. 1421--1441, 2019.
\bibitem{zhou2021} H. Zhou, S. Zhang, J. Peng, S. Zhang, J. Li, H. Xiong, and W. Zhang, ``Informer: Beyond efficient transformer for long sequence time-series forecasting,'' in \textit{Proceedings of the AAAI Conference on Artificial Intelligence}, vol. 35, no. 12, 2021, pp. 11106--11115.
\bibitem{wu2021} H. Wu, J. Xu, J. Wang, and M. Long, ``Autoformer: Decomposition transformers with auto-correlation for long-term series forecasting,'' \textit{Advances in Neural Information Processing Systems}, vol. 34, pp. 22419--22430, 2021.
\bibitem{lim2021} B. Lim and S. Zohren, ``Time-series forecasting with deep learning: a survey,'' \textit{Philosophical Transactions of the Royal Society A}, vol. 379, no. 2194, p. 20200209, 2021.
\bibitem{taylor2003} J. W. Taylor, ``Short-term electricity demand forecasting using double seasonal exponential smoothing,'' \textit{Journal of the Operational Research Society}, vol. 54, no. 8, pp. 799--805, 2003.
\bibitem{daut2017} M. A. M. Daut, M. Y. Hassan, H. Abdullah, H. A. Rahman, M. S. Majid, and A. H. Mohammad, ``Building electrical energy consumption forecasting analysis using conventional and machine learning methods,'' \textit{Journal of Electrical Systems and Information Technology}, vol. 4, no. 1, pp. 177--188, 2017.
\bibitem{mbamalu1993} G. A. N. Mbamalu and M. E. El-Hawary, ``Load forecasting via local linear regression with adaptive training window,'' \textit{IEEE Transactions on Power Systems}, vol. 8, no. 4, pp. 1662--1670, 1993.
\bibitem{drezga1998} I. Drezga and S. Rahman, ``Input variable selection for ANN-based short-term load forecasting,'' \textit{IEEE Transactions on Power Systems}, vol. 13, no. 4, pp. 1238--1244, 1998.
\end{thebibliography}

\appendices

\section{Failed Hypothesis Details}

All experiments in this appendix use the 6-fold expanding-window cross-validation framework described in Section IV-B. Results are reported as mean MAE across folds unless otherwise noted.

\subsection{Architecture Experiments (H1--H4)}
Each architecture was trained on the expanding-window setup. Hyperparameters were tuned on the Fold 1 validation set and fixed across folds. Training was performed using Adam (lr $= 10^{-3}$, batch size $= 64$, 100 epochs with early stopping).

\begin{table}[htbp]
\caption{Architecture Hypothesis Comparison (H1--H4)}
\label{tab:arch_failures}
\centering
\begin{tabular}{ccccc}
\toprule
\textbf{Hypothesis} & \textbf{Model} & \textbf{Mean MAE (MW)} & \textbf{Train Time} & \textbf{Params} \\
\midrule
H1 & DLinear (Baseline) & 93.9 & 12 mins & 42K \\
H1 & NLinear & 98.4 & 14 mins & 44K \\
H2 & DeepAR & 106.2 & 87 mins & 128K \\
H3 & DLinear + MLP & 94.1 & 22 mins & 126K \\
H4 & Transformer & 101.5 & 156 mins & 256K \\
\bottomrule
\end{tabular}
\end{table}

\begin{table}[htbp]
\caption{Per-Fold Breakdown (DLinear vs. Transformer)}
\label{tab:fold_breakdown_fail}
\centering
\begin{tabular}{cccc}
\toprule
\textbf{Fold} & \textbf{DLinear MAE (MW)} & \textbf{Transformer MAE (MW)} & \textbf{Difference} \\
\midrule
Fold 1 & 84.3 & 89.1 & +5.7\% \\
Fold 2 & 83.2 & 88.6 & +6.5\% \\
Fold 3 & 88.1 & 95.4 & +8.3\% \\
Fold 4 & 93.8 & 103.2 & +10.0\% \\
Fold 5 & 96.0 & 107.8 & +12.3\% \\
Fold 6 & 100.3 & 114.9 & +14.6\% \\
\bottomrule
\end{tabular}
\end{table}

The performance gap between DLinear and the Transformer widens in later folds (Table~\ref{tab:fold_breakdown_fail}). This is because high-capacity architectures overfit to historical patterns and struggle to generalize when the base load level shifts out-of-distribution.

\subsection{Feature Experiments (H5--H7)}
We incrementally add calendar features and weather features to the DLinear baseline (Table~\ref{tab:feature_failures}).

\begin{table}[htbp]
\caption{Feature Ablation Study}
\label{tab:feature_failures}
\centering
\begin{tabular}{cccc}
\toprule
\textbf{Feature Set} & \textbf{Mean MAE} & \textbf{vs Baseline} & \textbf{Folds Improved} \\
\midrule
DLinear only & 93.9 & --- & --- \\
+ Calendar (H5) & 86.4 & -8.0\% & 6 / 6 folds \\
+ Temperature (H7) & 89.2 & -5.0\% & 6 / 6 folds \\
+ Humidity & 92.8 & -1.2\% & 4 / 6 folds \\
+ Cloud Cover & 93.5 & -0.4\% & 3 / 6 folds \\
+ All Weather (H6) & 88.3 & -6.0\% & 4 / 6 folds \\
+ Calendar + All Weather & 84.7 & -9.8\% & 6 / 6 folds \\
\bottomrule
\end{tabular}
\end{table}

Adding temperature reduces the MAE by 5.0\%, which is smaller than in temperate climates, confirming weaker weather sensitivity in tropical regions \cite{mvungi2021}.

\subsection{Advanced Methods (H11--H12)}
We evaluate zero-shot foundation models (Chronos) and continual weight updates (EWC/SGD) (Table~\ref{tab:advanced_failures}).

\begin{table}[htbp]
\caption{Advanced Methods Comparison}
\label{tab:advanced_failures}
\centering
\begin{tabular}{ccc}
\toprule
\textbf{Model / Approach} & \textbf{Mean MAE (MW)} & \textbf{Operational Notes} \\
\midrule
Seasonal Naive (Ref) & 145.6 & No training \\
Chronos (Base) & 165.0 & Out-of-distribution failure \\
DLinear (Trained Baseline) & 93.9 & In-domain training \\
Online SGD (Daily) & Diverged & Diverged within 3 days \\
EWC ($\lambda = 100$) & 118.0 & Drift learning restricted \\
\bottomrule
\end{tabular}
\end{table}

Chronos captured diurnal cycles but failed to scale predictions, systematically underestimating the load scale. Daily online gradient updates diverged due to telemetry noise, while EWC constrained weight updates so severely that the model could not follow the growth trend.

\end{document}
